% TOP_PROBLEM: {"id":30007642,"problem_number":"LOCAL-30007642","title":"Reconciling income-inequality measurement","category_id":21,"category":{"id":21,"name":"economics","display_name":"Economics","description":"Open economic theory statements, published research agendas, and proposed scoped specifications; question provenance and historical priority are qualified per record.","slug":"economics","order_index":21,"created_at":"2026-10-08T00:00:00Z"},"status":"research_question","external_url":null,"legacy_ids":[],"published":false,"statement":"In the explicitly specified setting: How can surveys, tax records, and national accounts yield comparable inequality estimates when capital income, tax evasion, government services, and household sharing are treated differently? The mathematical statement above fixes the target, assumptions and scope of this catalogue formulation.","economics":{"original_id":131,"field":"Inequality and economic mobility","kind":"identification","formalization_basis":"adapted_research_question","source_status":"research_agenda","priority_status":"earliest_found","priority_note":"Earliest directly inspected qualifying passage in this audit; earlier origins were not established. A review or research-agenda statement does not imply original priority.","earliest_verified_reference":{"authors":["Conor Clarke","Wojciech Kopczuk"],"title":"Measuring Income and Income Inequality","year":2025,"url":"https://repository.library.washu.edu/law_scholarship/812/","doi":"10.1257/jep.20241424","locator":"Abstract, paragraphs 1–2; published article pp. 103–126","evidence_summary":"The authors identify disputed concepts and data-source tradeoffs, while rejecting a uniquely comprehensive, controversy-free income measure."},"current_reference":{"authors":["Conor Clarke","Wojciech Kopczuk"],"title":"Measuring Income and Income Inequality","year":2025,"url":"https://repository.library.washu.edu/law_scholarship/812/","doi":"10.1257/jep.20241424","locator":"Abstract, paragraphs 1–2; published article pp. 103–126","evidence_summary":"The authors identify disputed concepts and data-source tradeoffs, while rejecting a uniquely comprehensive, controversy-free income measure."},"formalization_note":"The cited passage establishes a research direction, not this exact mathematical statement. The notation, population, policy menu, assumptions, and estimands are an original adaptation for this catalogue; no universal unsolved theorem or source priority is claimed."},"statusQualification":"Published question or research agenda verified at the cited locator. The mathematical specification is adapted for this catalogue and includes additional assumptions; source publication does not establish first-ever priority or certify this exact model remains unsolved. The cited passage establishes a research direction, not this exact mathematical statement. The notation, population, policy menu, assumptions, and estimands are an original adaptation for this catalogue; no universal unsolved theorem or source priority is claimed.","statusReviewedAt":"2026-10-08"}
% FIRST_POSTED: 2026-10-08
% ENTRY_KIND: statement_only
% !TeX program = lualatex
\documentclass[11pt]{article}
\usepackage{fontspec}
\usepackage[a4paper,margin=25mm]{geometry}
\usepackage{amsmath,amssymb,mathtools}
\usepackage{xurl}
\usepackage[hidelinks]{hyperref}
\newcommand{\sref}[2]{\hyperlink{source:#1:#2}{[S#2]}}
\setlength{\emergencystretch}{3em}
\begin{document}
% problemId:30007642
\section*{Reconciling income-inequality measurement}\label{30007642}
\subsection{Definitions and mathematical statement}
Fix a country, accounting period, adult population, and a declared income concept \(\theta\). The concept specifies capital gains, retained business earnings, imputed government services, household sharing, taxes, and transfers. Let \(y_i(\theta)\) be person \(i\)'s income under those definitions, and let \(F_\theta\) be its population distribution. Observed sources comprise tax records \(D_T\), household surveys \(D_S\), and aggregate national accounts \(D_N\), each with stated coverage and error processes. For \(p\in(0,1)\), define the top-share functional \(Q_p(F)=\int_{1-p}^{1}F^{-1}(u)\,du/\int_0^1F^{-1}(u)\,du\), assuming a positive finite mean. The measurement task is to identify or bound \(Q_p(F_\theta)\) and reconcile discrepancies across sources without treating an accounting allocation convention as observed household income. Quantify sensitivity to nonfilers, evasion, survey nonresponse, asset valuation, and income-sharing assumptions. Separate disagreements about \(\theta\) from statistical uncertainty conditional on \(\theta\). Construct a transparent set of estimates for a predeclared family of concepts. There need not exist one uniquely correct comprehensive measure; the research problem is reproducible measurement and defensible reconciliation, with remaining normative choices made explicit.

\par\textbf{Formulation and scope.} The cited passage establishes a research direction, not this exact mathematical statement. The notation, population, policy menu, assumptions, and estimands are an original adaptation for this catalogue; no universal unsolved theorem or source priority is claimed.

\par\textbf{Open-question provenance.} An explicit question or research agenda was verified at the source locator. This does not by itself certify that every later version remains unresolved \sref{30007642}{1}.

\par\textbf{Historical priority.} Earliest directly inspected qualifying passage in this audit; earlier origins were not established. A review or research-agenda statement does not imply original priority.

\subsection{Short English statement}
In the explicitly specified setting: How can surveys, tax records, and national accounts yield comparable inequality estimates when capital income, tax evasion, government services, and household sharing are treated differently? The mathematical statement above fixes the target, assumptions and scope of this catalogue formulation.

\subsection{Sources}
\begin{itemize}\raggedright
\item[\textnormal{[S1]}]\hypertarget{source:30007642:1}{} Conor Clarke, Wojciech Kopczuk. 2025. Measuring Income and Income Inequality. Abstract, paragraphs 1–2; published article pp. 103–126.
\url{https://repository.library.washu.edu/law_scholarship/812/}.
DOI: 10.1257/jep.20241424.
{\small\textit{Earliest verified question/agenda reference; historical priority qualified below. The authors identify disputed concepts and data-source tradeoffs, while rejecting a uniquely comprehensive, controversy-free income measure.}}
\end{itemize}
\end{document}
